Modern genomics has enabled large-scale measurements of genetic variation and cellular dynamics, offering unprecedented opportunities to study biological systems. However, many scientific questions require forecasting beyond observed data, such as predicting the number of undiscovered genetic variants in future studies or inferring and extrapolating cellular dynamics from snapshot level measurements. These problems are fundamentally challenging due to limited observations, high noise, and the partial observability of underlying processes. At the same time, domain knowledge is often available but incomplete or imperfect, raising the need for methods that can flexibly integrate such information.

In this thesis, we develop probabilistic frameworks for forecasting under limited and indirect observations, with a focus on two central problems in genomics. First, we propose a predictor of novel genetic variants in follow-up sequencing studies across multiple populations. We propose Bayesian nonparametric models that capture both shared and population-specific variation while avoiding restrictive assumptions on the total number of variants. Our approach incorporates domain knowledge such as power-law-like growth patterns while allowing data-driven adaptation through an empirical Bayes-like strategy.

Second, we address the problem of inferring and forecasting latent stochastic dynamics from snapshot observations, as arises in single-cell sequencing. Since individual cells cannot be tracked over time, learning the underlying dynamics is inherently ill-posed. We develop methods based on flexible stochastic differential equation models that integrate domain knowledge without requiring it to be fully specified. Our approach enables reconstruction of trajectory distributions and forecasting of their future evolution.

Across these settings, the key methodological challenge is to balance flexibility with structure: models must be expressive enough to fit complex data while remaining sufficiently constrained to incorporate imperfect prior knowledge. By developing approaches that operate in this regime, this thesis advances a general framework for reliable forecasting in scientific domains characterized by high noise and partial observability.